\documentclass[10pt,a4paper]{article}
\usepackage[margin=1.8cm]{geometry}
\usepackage{fontspec,xeCJK}
\usepackage{amsmath,amssymb,booktabs,caption}
\usepackage[hidelinks]{hyperref}
\setCJKmainfont{Songti TC}
\setCJKsansfont{Heiti TC}
\setCJKmonofont{Heiti TC}
\setlength{\parindent}{0pt}
\setlength{\parskip}{0.4em}
\renewcommand{\arraystretch}{1.13}
\captionsetup{font=small,skip=5pt}
\setlength{\abovedisplayskip}{6pt}
\setlength{\belowdisplayskip}{6pt}

\begin{document}
\begin{center}
{\Large\bfseries HW6 AdaBoost 計算說明}\par
固定抽樣與原始樣本加權誤差
\end{center}
本報告依學習單完成三輪 AdaBoost 計算，選出的分類器依序為 A、E、D。

四筆新資料的預測依序為 $-1,+1,-1,+1$。
以下保留資料、選擇過程及權重更新，供逐項核對。

\section*{一、原始訓練資料}
\begin{center}
\captionof{table}{題目指定的十筆訓練資料}\label{tab:training}
\begin{tabular}{rrrrrrrrrrr}\toprule
ID & 1&2&3&4&5&6&7&8&9&10\\
$x$ & 1&2&3&4&5&6&7&8&9&10\\
$y$ & $-1$&$-1$&$-1$&$-1$&$+1$&$+1$&$+1$&$-1$&$-1$&$+1$\\
\bottomrule\end{tabular}
\end{center}
表\ref{tab:training} 的類別隨 $x$ 有多次轉換，單一 stump 無法將全部樣本分對。

\section*{二、候選規則與計算方式}
只使用切分點 4.5、7.5、9.5，各搭配兩種左右側標籤，共 A--F 六個規則，
詳見各輪候選規則表。每輪先在固定抽樣集比較未加權錯誤筆數；
若同一 ID 抽到多次，每次都計入。選定 $h_t$ 後，才在原始十筆資料計算
\begin{align*}
\epsilon_t&=\sum_{i=1}^{10}w_i^{(t)}\,I\{y_i\ne h_t(x_i)\}, &
\alpha_t&=\frac12\ln\frac{1-\epsilon_t}{\epsilon_t},\\
u_i&=w_i^{(t)}\exp[-\alpha_t y_i h_t(x_i)], &
w_i^{(t+1)}&=\frac{u_i}{\sum_{j=1}^{10}u_j}.
\end{align*}
如此可保留完整的權重分布，不以有限次抽樣的錯誤比例取代加權錯誤率。
本題各輪均有唯一最佳規則；程式在候選同分時依 A--F 順序選第一個，
投票若恰為零則回傳 0，此兩項慣例皆未影響本例結果。

\textbf{精度與重現：}中間使用 Python 雙精度浮點數，表格顯示六位小數；
公式中的分數用於核對精確值。固定抽樣清單由題目給定，不重新抽樣，
因此結果不依賴亂數種子。資料、候選及結果儲存於 \texttt{results/adaboost.json}。
\clearpage
\section*{三、Round 1}
初始權重皆為 $0.1$。本輪抽樣 ID 為 \[[1,\,2,\,3,\,4,\,5,\,6,\,7,\,8,\,9,\,10].\]
\begin{center}
\captionof{table}{Round 1 的候選規則比較}
\label{tab:round1}
\begin{tabular}{rrrrr}
\toprule
規則 & 切分點 & 左側預測 & 右側預測 & 抽樣錯誤筆數 \\
\midrule
A & 4.5 & $-1$ & $+1$ & 2 \\
B & 4.5 & $+1$ & $-1$ & 8 \\
C & 7.5 & $-1$ & $+1$ & 5 \\
D & 7.5 & $+1$ & $-1$ & 5 \\
E & 9.5 & $-1$ & $+1$ & 3 \\
F & 9.5 & $+1$ & $-1$ & 7 \\
\bottomrule
\end{tabular}
\end{center}
最少錯誤為 2 筆，選出規則 A，其分類器為
\[h_1(x)=\begin{cases}-1,&x\le 4.5,\\+1,&x>4.5.\end{cases}\]
回到原始 10 筆資料，分錯 ID 為 8, 9，故
\[\epsilon_1=2(0.1)=\frac15\approx 0.200000,\qquad \alpha_1=\frac12\ln\left(4\right)\approx 0.693147.\]
正規化常數 $Z_1=\sum_i u_i=0.800000$。以分錯的 ID 8 為例，$y_i h_1(x_i)=-1$，因此
\[u_{8}=0.100000\exp(+0.693147)\approx 0.200000,\qquad w_{8}^{(2)}=\frac{u_{8}}{Z_1}\approx 0.250000.\]
分對樣本的新權重均為 $1/16=0.062500$；分錯樣本（ID 8、9）的新權重均為 $1/4=0.250000$。
\begin{center}
\captionof{table}{Round 1 在原始資料上的權重更新}
\label{tab:weights1}
\begin{tabular}{rrrrrr}
\toprule
ID & $y_i$ & $h_1(x_i)$ & $w_i^{(1)}$ & $u_i$ & $w_i^{(2)}$ \\
\midrule
1 & $-1$ & $-1$ & 0.100000 & 0.050000 & 0.062500 \\
2 & $-1$ & $-1$ & 0.100000 & 0.050000 & 0.062500 \\
3 & $-1$ & $-1$ & 0.100000 & 0.050000 & 0.062500 \\
4 & $-1$ & $-1$ & 0.100000 & 0.050000 & 0.062500 \\
5 & $+1$ & $+1$ & 0.100000 & 0.050000 & 0.062500 \\
6 & $+1$ & $+1$ & 0.100000 & 0.050000 & 0.062500 \\
7 & $+1$ & $+1$ & 0.100000 & 0.050000 & 0.062500 \\
8 & $-1$ & $+1$ & 0.100000 & 0.200000 & 0.250000 \\
9 & $-1$ & $+1$ & 0.100000 & 0.200000 & 0.250000 \\
10 & $+1$ & $+1$ & 0.100000 & 0.050000 & 0.062500 \\
\bottomrule
\end{tabular}
\end{center}
權重總和為 $8(1/16)+2(1/4)=1$；兩筆錯分樣本合計占一半抽樣機率。

\clearpage
\section*{四、Round 2}
使用 Round 1 更新後的權重。固定抽樣 ID 為 \[[1,\,3,\,4,\,5,\,8,\,8,\,9,\,9,\,9,\,10].\]
\begin{center}
\captionof{table}{Round 2 的候選規則比較}
\label{tab:round2}
\begin{tabular}{rrrrr}
\toprule
規則 & 切分點 & 左側預測 & 右側預測 & 抽樣錯誤筆數 \\
\midrule
A & 4.5 & $-1$ & $+1$ & 5 \\
B & 4.5 & $+1$ & $-1$ & 5 \\
C & 7.5 & $-1$ & $+1$ & 6 \\
D & 7.5 & $+1$ & $-1$ & 4 \\
E & 9.5 & $-1$ & $+1$ & 1 \\
F & 9.5 & $+1$ & $-1$ & 9 \\
\bottomrule
\end{tabular}
\end{center}
最少錯誤為 1 筆，選出規則 E，其分類器為
\[h_2(x)=\begin{cases}-1,&x\le 9.5,\\+1,&x>9.5.\end{cases}\]
回到原始 10 筆資料，分錯 ID 為 5, 6, 7，故
\[\epsilon_2=3\left(\frac1{16}\right)=\frac3{16}\approx 0.187500,\qquad \alpha_2=\frac12\ln\left(\frac{13}{3}\right)\approx 0.733169.\]
正規化常數 $Z_2=\sum_i u_i=0.780625$。以分錯的 ID 5 為例，$y_i h_2(x_i)=-1$，因此
\[u_{5}=0.062500\exp(+0.733169)\approx 0.130104,\qquad w_{5}^{(3)}=\frac{u_{5}}{Z_2}\approx 0.166667.\]
\begin{center}
\captionof{table}{Round 2 在原始資料上的權重更新}
\label{tab:weights2}
\begin{tabular}{rrrrrr}
\toprule
ID & $y_i$ & $h_2(x_i)$ & $w_i^{(2)}$ & $u_i$ & $w_i^{(3)}$ \\
\midrule
1 & $-1$ & $-1$ & 0.062500 & 0.030024 & 0.038462 \\
2 & $-1$ & $-1$ & 0.062500 & 0.030024 & 0.038462 \\
3 & $-1$ & $-1$ & 0.062500 & 0.030024 & 0.038462 \\
4 & $-1$ & $-1$ & 0.062500 & 0.030024 & 0.038462 \\
5 & $+1$ & $-1$ & 0.062500 & 0.130104 & 0.166667 \\
6 & $+1$ & $-1$ & 0.062500 & 0.130104 & 0.166667 \\
7 & $+1$ & $-1$ & 0.062500 & 0.130104 & 0.166667 \\
8 & $-1$ & $-1$ & 0.250000 & 0.120096 & 0.153846 \\
9 & $-1$ & $-1$ & 0.250000 & 0.120096 & 0.153846 \\
10 & $+1$ & $+1$ & 0.062500 & 0.030024 & 0.038462 \\
\bottomrule
\end{tabular}
\end{center}
更新後 ID 1--4、10 各為 $1/26$，ID 5--7 各為 $1/6$，ID 8、9 各為 $2/13$。因此總和為
\[5\left(\frac1{26}\right)+3\left(\frac16\right)+2\left(\frac2{13}\right)=1.\]
注意：抽樣集只有 ID 5 的一筆錯誤（$1/10$）；原始資料上 ID 5、6、7 的權重合計才是 $\epsilon_2=3/16$。

\clearpage
\section*{五、Round 3}
使用 Round 2 更新後的權重。固定抽樣 ID 為 \[[1,\,5,\,5,\,6,\,6,\,7,\,7,\,8,\,9,\,9].\]
\begin{center}
\captionof{table}{Round 3 的候選規則比較}
\label{tab:round3}
\begin{tabular}{rrrrr}
\toprule
規則 & 切分點 & 左側預測 & 右側預測 & 抽樣錯誤筆數 \\
\midrule
A & 4.5 & $-1$ & $+1$ & 3 \\
B & 4.5 & $+1$ & $-1$ & 7 \\
C & 7.5 & $-1$ & $+1$ & 9 \\
D & 7.5 & $+1$ & $-1$ & 1 \\
E & 9.5 & $-1$ & $+1$ & 6 \\
F & 9.5 & $+1$ & $-1$ & 4 \\
\bottomrule
\end{tabular}
\end{center}
最少錯誤為 1 筆，選出規則 D，其分類器為
\[h_3(x)=\begin{cases}+1,&x\le 7.5,\\-1,&x>7.5.\end{cases}\]
回到原始 10 筆資料，分錯 ID 為 1, 2, 3, 4, 10，故
\[\epsilon_3=5\left(\frac1{26}\right)=\frac5{26}\approx 0.192308,\qquad \alpha_3=\frac12\ln\left(\frac{21}{5}\right)\approx 0.717542.\]
規則 D 對原始資料分錯 ID 1、2、3、4、10，共 \textbf{5 筆}；抽樣錯誤只有 1 筆，兩者的評估資料不同。本題不再更新第三輪權重。

\section*{六、三個分類器的加權投票}
令 $S(x)=\sum_{t=1}^3\alpha_t h_t(x)$、$H(x)=\operatorname{sign}(S(x))$。
\begin{center}
\captionof{table}{新資料的完整預測結果}
\label{tab:predictions}
\begin{tabular}{rrrrrr}
\toprule
$x$ & $h_1(x)$ & $h_2(x)$ & $h_3(x)$ & $S(x)$ & $H(x)$ \\
\midrule
2.5 & $-1$ & $-1$ & $+1$ & -0.708773 & $-1$ \\
5.5 & $+1$ & $-1$ & $+1$ & 0.677521 & $+1$ \\
8.5 & $+1$ & $-1$ & $-1$ & -0.757564 & $-1$ \\
9.8 & $+1$ & $+1$ & $-1$ & 0.708773 & $+1$ \\
\bottomrule
\end{tabular}
\end{center}
例如 $x=5.5$，三個分類器的預測依序為 $+1,-1,+1$，因此
\begin{align*}S(5.5)&=\tfrac12\ln(4)-\tfrac12\ln(13/3)+\tfrac12\ln(21/5)\\&\approx 0.693147-0.733169+0.717542=0.677521>0,\qquad H(5.5)=+1.\end{align*}
四筆資料的加權分數皆不為零，因此預測沒有平手情形。

\clearpage
\section*{七、思考題}
\textbf{Q1. 第一輪分錯的樣本，為何第二輪容易重複抽到？}

ID 8、9 的權重由 0.1 升至 0.25，每筆是分對樣本權重 0.0625 的四倍；若依此機率有放回抽樣，就較容易重複抽到它們。
本例第二輪的固定清單中，ID 8 出現兩次、ID 9 出現三次，但權重較高不保證每次抽樣都會重複。

\textbf{Q2. 為何計算 $\epsilon_2$ 仍要回到原始資料及其權重？}

抽樣資料只用來選規則，$\epsilon_2$ 的定義則是原始 10 筆資料中錯分樣本的權重總和。
本例規則 E 的抽樣錯誤率為 $1/10$，但原始資料錯分 ID 5、6、7，故 $\epsilon_2=3/16$。

\textbf{Q3. 為何第二輪可能選出不同切分規則？}

抽樣組成改變後，各規則的未加權錯誤筆數也隨之改變；第一輪 A 錯 2 筆、E 錯 3 筆，所以選 A。
第二輪重複抽到 ID 8、9，使 A 錯 5 筆，而 E 只錯 1 筆，因此改選切分點 9.5 的 E。

\section*{八、重現與驗證}
在 HW6 目錄執行 \texttt{uv run python src/adaboost.py} 即可重建 CSV、JSON、
console 摘要及報告的計算分節，再於 \texttt{report/} 執行 tectonic 編譯兩份 PDF。
完整指令與驗收對照見 \texttt{README.md} 及 \texttt{docs/verification.md}。

自動化測試以手算分數核對三輪加權誤差、分類器權重與樣本權重，
並檢查重複抽樣、切分點等號、權重傳遞及四筆投票。
PDF 的可讀性與版面另以實際渲染圖片檢查；自動化數值測試不代表版面已驗證。

\textbf{工具使用說明：}本作業使用 OpenAI Codex 協助撰寫 Python 計算、
核對結果及編排 LaTeX；原始學習單保留不變，未使用外部資料或隨機實驗。
\end{document}
